%%%%%%%%%%%%%%%%%%%%%%%%%%%%%%%%%%%%%%%%%
% Beamer Presentation
% LaTeX Template
% Version 1.0 (10/11/12)
%
% This template has been downloaded from:
% http://www.LaTeXTemplates.com
%
% License:
% CC BY-NC-SA 3.0 (http://creativecommons.org/licenses/by-nc-sa/3.0/)
%
%%%%%%%%%%%%%%%%%%%%%%%%%%%%%%%%%%%%%%%%%

%----------------------------------------------------------------------------------------
%	PACKAGES AND THEMES
%----------------------------------------------------------------------------------------

\documentclass[xcolor=dvipsnames]{beamer}

\mode<presentation> {

% The Beamer class comes with a number of default slide themes
% which change the colors and layouts of slides. Below this is a list
% of all the themes, uncomment each in turn to see what they look like.

%\usetheme{default}
%\usetheme{AnnArbor}
%\usetheme{Antibes}
%\usetheme{Bergen}
%\usetheme{Berkeley}
%\usetheme{Berlin}
%\usetheme{Boadilla}
%\usetheme{CambridgeUS}
%\usetheme{Copenhagen}
%\usetheme{Darmstadt}
%\usetheme{Dresden}
%\usetheme{Frankfurt}
%\usetheme{Goettingen}
%\usetheme{Hannover}
%\usetheme{Ilmenau}
%\usetheme{JuanLesPins}
%\usetheme{Luebeck}
\usetheme{Madrid}
%\usetheme{Malmoe}
%\usetheme{Marburg}
%\usetheme{Montpellier}
%\usetheme{PaloAlto}
%\usetheme{Pittsburgh}
%\usetheme{Rochester}
%\usetheme{Singapore}
%\usetheme{Szeged}
%\usetheme{Warsaw}

% As well as themes, the Beamer class has a number of color themes
% for any slide theme. Uncomment each of these in turn to see how it
% changes the colors of your current slide theme.

%\usecolortheme{albatross}
%\usecolortheme{beaver}
%\usecolortheme{beetle}
%\usecolortheme{crane}
\usecolortheme{dolphin}
%\usecolortheme{dove}
%\usecolortheme{fly}
%\usecolortheme{lily}
%\usecolortheme{orchid}
%\usecolortheme{rose}
%\usecolortheme{seagull}
%\usecolortheme{seahorse}
%\usecolortheme{whale}
%\usecolortheme{wolverine}

%\setbeamertemplate{footline} % To remove the footer line in all slides uncomment this line
%\setbeamertemplate{footline}[page number] % To replace the footer line in all slides with a simple slide count uncomment this line

\setbeamertemplate{navigation symbols}{} % To remove the navigation symbols from the bottom of all slides uncomment this line
}



\input inputs/packages.tex
\input inputs/abbreviations.tex
\input inputs/commands.tex

%----------------------------------------------------------------------------------------
%	TITLE PAGE
%----------------------------------------------------------------------------------------

\newcommand\meeting{Thesis defense}

\title[\meeting]{Measurement of the photon energy spectrum in
inclusive radiative \texorpdfstring{$B$}{B} meson decays using the
hadronic-tagging method} % The short title appears at the bottom of every slide, the full title is only on the title page

\author[Henrikas Svidras]{\texorpdfstring{\footnotesize \underline{Henrikas Svidras}}{Henrikas Svidras}} % Your name
%\subject{\text{\small KEK, Tsukuba}}
\institute[DESY] % Your institution as it will appear on the bottom of every slide, may be shorthand to save space
{
\text{\large \meeting}
\vspace{-10pt}
}
\titlegraphic{\includegraphics[width=2cm]{logos/DESY.png}\hspace*{4.75cm}~%
   \includegraphics[width=2cm]{logos/belle2.png}
}

\date{March 13, 2023} % Date, can be changed to a custom date

\begin{document}

\setbeamercolor{background canvas}{bg=}
{\setbeamertemplate{footline}{} 
\begin{frame}
\titlepage % Print the title page as the first slide
\end{frame}
}
\addtocounter{framenumber}{-1}

%----------------------------------------------------------------------------------------
%	PRESENTATION SLIDES
%----------------------------------------------------------------------------------------
\section{Introduction}
\begin{frame}{Introduction: radiative \safeB meson decays}
\centering
\textbf{What are radiative $\bm{B}$ decays?}

   \begin{tikzpicture}
    \begin{feynman}
    \vertex (i1){b};
    \vertex[right =1.2cm of i1] (a);
    \vertex[right=0.85cm of a] (b);
    \vertex[right=0.85cm of b] (c);
    \vertex[right=1.2cm of c] (o1) {s,d};
    \vertex[below=2em of c] (g1);
    \vertex[right=2em of g1] (o2) {$\gamma$};
    
    \diagram* {
    (i1) -- [fermion] (a) -- [fermion, edge label =\(uct\)] (c) -- [fermion] (o1),
    (a) -- [boson,half left, edge label = \(W^{\pm}\)] (c),
    (b) -- [photon] (o2),
    };
    \end{feynman}
\end{tikzpicture}
    
   \begin{tikzpicture}
    \begin{feynman}
    \vertex (i1){b};
    \vertex[right =1.2cm of i1] (a) ;
    \vertex[right=0.85cm of a] (b);
    \vertex[right=0.85cm of b] (c);
    \vertex[right=1.2cm of c] (o1) {s,d};
    \vertex[below=2em of c] (g1);
    \vertex[right=2em of g1] (o2) {$\gamma$};
    
    \diagram* {
    (i1) -- [fermion] (a) -- [boson, edge label =\(W^{\pm}\)] (c) -- [fermion] (o1),
    (a) -- [fermion,half left, edge label = \(uct\)] (c),
    (b) -- [photon] (o2),
    };
    \end{feynman}
\end{tikzpicture}
    
   \begin{tikzpicture}
    \begin{feynman}
    \vertex (i1){b};
    \vertex[right =1.2cm of i1] (a) ;
    \vertex[right=0.85cm of a] (b);
    \vertex[right=0.85cm of b] (c);
    \vertex[right=1.2cm of c] (o1) {s,d};
    \vertex[below=2em of c] (g1);
    \vertex[right=2em of g1] (o2) {$\gamma$};
    
    \diagram* {
    (i1) -- [fermion] (a) -- [scalar, edge label =\(H^{\pm}\)] (c) -- [fermion] (o1),
    (a) -- [fermion,half left, edge label = \(uct\)] (c),
    (b) -- [photon] (o2),
    };
    \end{feynman}
\end{tikzpicture}
\end{frame}

\end{document}